% 图 56.1 —— 由 `checks/emit_hand.py` 自动拼出（probe 精测 + prims2 图元分解）
%%
% ⚠ 这是**稿子不是成品**：跑 `checks/checkhand.py 56.1` 看指标 + 看叠合图。
%%
%% ⚠ 待核：
%   · 本图 probe 报出阴影族 1 个 —— **本脚本不生成阴影**（要照原书手写）；prims2 的 strokes 里可能含阴影线，核对时留意别画重
%   · 可疑字形 2 项（空心圈 ⊙ 常被认成字母 o）—— 见文件末
%%
\documentclass[12pt]{standalone}
\usepackage{fontspec}
\setsansfont{Arial}
\newfontfamily\mfallback{DejaVu Sans}
\usepackage{tikz}
\begin{document}
\begin{tikzpicture}[x=1pt, y=-1pt, line width=5.6pt, line cap=round, line join=round]

  %% ════ 填充区（prims2 的 fills）════

  %% ════ 直线段（prims2 的 strokes）════
  \draw[line width=4.0pt] (708.0,206.0) -- (714.0,204.0);
  \draw[line width=5.0pt] (1009.0,209.0) -- (1015.0,205.0);
  \draw[line width=7.0pt] (813.0,106.0) -- (822.0,101.0);
  \draw[line width=8.1pt] (973.0,301.0) -- (980.0,304.0);
  \draw[line width=7.0pt] (754.0,142.0) -- (745.0,146.0);
  \draw[line width=5.0pt] (640.0,323.0) -- (625.0,328.0);
  \draw[line width=6.0pt] (658.0,178.0) -- (673.0,185.0);
  \draw[line width=4.5pt] (747.0,129.0) -- (747.0,136.0);
  \draw[line width=7.0pt] (1015.0,204.0) -- (1015.0,197.0);
  \draw[line width=7.0pt] (1031.0,243.0) -- (1029.0,235.0);
  \draw[line width=6.7pt] (1057.0,114.0) -- (1051.0,116.0);
  \draw[line width=4.5pt] (760.0,141.0) -- (755.0,142.0);
  \draw[line width=3.0pt] (730.0,235.0) -- (730.0,232.0);
  \draw[line width=5.0pt] (729.0,236.0) -- (692.0,259.0);
  \draw[line width=7.0pt] (1001.0,370.8) -- (991.2,382.8);
  \draw[line width=7.0pt] (991.2,382.8) -- (989.0,390.0);
  \draw[line width=6.0pt] (963.0,135.0) -- (958.0,128.0);
  \draw[line width=5.0pt] (747.0,280.0) -- (727.0,374.0);
  \draw[line width=7.0pt] (754.0,136.0) -- (778.0,119.0);
  \draw[line width=7.0pt] (1016.0,205.0) -- (1020.0,216.0);
  \draw[line width=4.0pt] (622.0,328.0) -- (613.0,333.0);
  \draw[line width=4.0pt] (1058.0,115.0) -- (1052.0,142.0);
  \draw[line width=5.0pt] (753.0,137.0) -- (748.0,137.0);
  \draw[line width=6.7pt] (338.0,174.0) -- (344.0,174.0);
  \draw[line width=6.0pt] (707.0,207.0) -- (700.0,214.0);
  \draw[line width=10.8pt] (722.0,170.0) -- (719.0,161.0);
  \draw[line width=3.0pt] (754.0,138.0) -- (754.0,141.0);
  \draw[line width=6.0pt] (698.0,205.0) -- (699.0,214.0);
  \draw[line width=6.7pt] (699.0,204.0) -- (705.0,202.0);
  \draw[line width=5.0pt] (651.0,181.0) -- (656.0,178.0);
  \draw[line width=5.0pt] (791.0,155.0) -- (813.0,107.0);
  \draw[line width=6.5pt] (779.0,104.0) -- (778.0,117.0);
  \draw[line width=4.0pt] (696.0,199.0) -- (698.0,204.0);
  \draw[line width=5.0pt] (1019.0,217.0) -- (1009.0,213.0);
  \draw[line width=5.0pt] (747.0,270.0) -- (764.0,219.0);
  \draw[line width=4.0pt] (769.0,210.0) -- (785.0,163.0);
  \draw[line width=5.0pt] (856.0,444.0) -- (904.0,387.0);
  \draw[line width=4.0pt] (979.0,305.0) -- (910.0,385.0);
  \draw[line width=5.0pt] (846.0,447.0) -- (843.9,452.1);
  \draw[line width=8.7pt] (726.0,392.8) -- (729.0,385.0);
  \draw[line width=5.0pt] (940.0,392.0) -- (910.0,386.0);
  \draw[line width=6.6pt] (622.0,335.0) -- (624.0,329.0);
  \draw[line width=5.0pt] (981.0,303.0) -- (1022.9,251.1);
  \draw[line width=7.8pt] (1022.9,251.1) -- (1030.0,249.0);
  \draw[line width=7.0pt] (791.0,156.0) -- (785.0,162.0);
  \draw[line width=6.5pt] (1037.0,236.0) -- (1032.0,234.0);
  \draw[line width=5.0pt] (976.0,139.0) -- (1051.0,143.0);
  \draw[line width=5.0pt] (825.0,91.2) -- (822.0,100.0);
  \draw[line width=5.7pt] (981.0,305.0) -- (983.0,310.0);
  \draw[line width=6.8pt] (682.9,272.1) -- (689.0,270.0);
  \draw[line width=6.0pt] (786.0,117.0) -- (812.0,106.0);
  \draw[line width=4.7pt] (707.0,206.0) -- (705.0,202.0);
  \draw[line width=5.3pt] (657.0,176.0) -- (656.0,171.0);
  \draw[line width=5.0pt] (943.0,394.0) -- (977.0,396.0);
  \draw[line width=5.0pt] (763.0,218.0) -- (731.0,231.0);
  \draw[line width=5.0pt] (835.0,353.0) -- (830.0,353.0);
  \draw[line width=4.0pt] (905.0,386.0) -- (909.0,386.0);
  \draw[line width=6.3pt] (779.0,118.0) -- (785.0,117.0);
  \draw[line width=9.0pt] (1030.0,234.0) -- (1020.0,217.0);
  \draw[line width=7.7pt] (744.0,145.0) -- (746.0,138.0);
  \draw[line width=6.7pt] (947.0,181.0) -- (949.0,187.0);
  \draw[line width=7.9pt] (1032.0,243.0) -- (1037.0,237.0);
  \draw[line width=7.5pt] (988.0,391.0) -- (977.0,396.0);
  \draw[line width=9.7pt] (724.0,171.0) -- (733.0,171.0);
  \draw[line width=6.0pt] (747.0,271.0) -- (748.0,279.0);
  \draw[line width=6.0pt] (769.0,211.0) -- (764.0,217.0);

  %% ════ 曲线（prims2 的 curves —— 三段式，坐标同为 (y,x)）════
  \draw[line width=5.0pt] (730.0,230.0) .. controls (726.3,196.6) and (755.4,170.9) .. (784.0,162.0);
  \draw[line width=5.0pt] (904.0,386.0) .. controls (880.9,376.4) and (855.3,369.6) .. (836.0,353.0);
  \draw[line width=5.0pt] (338.0,178.0) .. controls (389.6,257.0) and (360.6,374.4) .. (278.0,422.0);
  \draw[line width=5.0pt] (278.0,422.0) .. controls (202.2,467.4) and (93.4,441.4) .. (48.0,362.0);
  \draw[line width=5.0pt] (48.0,362.0) .. controls (2.6,287.8) and (21.1,182.7) .. (92.6,129.4);
  \draw[line width=5.0pt] (92.6,129.4) .. controls (166.9,71.5) and (287.8,90.7) .. (337.0,174.0);
  \draw[line width=4.0pt] (829.0,354.0) .. controls (799.9,372.4) and (765.3,376.8) .. (733.0,385.0);
  \draw[line width=4.0pt] (948.0,180.0) .. controls (968.7,187.0) and (989.7,195.7) .. (1008.0,209.0);
  \draw[line width=5.0pt] (962.0,136.0) .. controls (903.9,137.9) and (845.6,139.7) .. (792.0,156.0);
  \draw[line width=4.0pt] (1039.0,237.0) .. controls (1056.5,256.7) and (1066.2,281.4) .. (1069.0,306.9);
  \draw[line width=5.0pt] (1069.0,306.9) .. controls (1065.1,337.5) and (1048.3,371.6) .. (1016.0,381.0);
  \draw[line width=4.0pt] (1016.0,381.0) .. controls (1008.4,386.9) and (999.6,390.4) .. (990.0,391.0);
  \draw[line width=8.0pt] (1031.0,249.0) .. controls (1042.9,293.1) and (1018.1,333.1) .. (1001.0,370.8);
  \draw[line width=5.0pt] (946.0,180.0) .. controls (884.8,165.8) and (825.2,191.1) .. (770.0,211.0);
  \draw[line width=6.0pt] (958.0,128.0) .. controls (919.3,100.9) and (870.1,91.5) .. (823.0,101.0);
  \draw[line width=4.0pt] (613.0,333.0) .. controls (551.9,348.2) and (481.8,347.5) .. (423.0,326.0);
  \draw[line width=4.0pt] (749.0,280.0) .. controls (771.8,306.9) and (800.0,332.6) .. (829.0,352.0);
  \draw[line width=4.0pt] (656.0,177.0) .. controls (584.7,150.6) and (498.4,149.9) .. (425.0,171.0);
  \draw[line width=4.0pt] (843.9,452.1) .. controls (820.5,474.5) and (791.3,499.1) .. (758.0,488.0);
  \draw[line width=4.0pt] (758.0,488.0) .. controls (724.9,468.8) and (725.6,427.6) .. (726.0,392.8);
  \draw[line width=4.0pt] (1058.0,113.0) .. controls (1060.2,91.3) and (1057.9,68.9) .. (1045.7,50.7);
  \draw[line width=4.0pt] (1045.7,50.7) .. controls (981.6,-13.8) and (871.4,20.7) .. (825.0,91.2);
  \draw[line width=5.0pt] (726.0,383.0) .. controls (700.3,386.0) and (667.1,381.9) .. (656.2,355.2);
  \draw[line width=4.0pt] (656.2,355.2) .. controls (647.5,324.5) and (665.3,295.7) .. (682.9,272.1);
  \draw[line width=4.0pt] (1039.0,236.0) .. controls (1056.9,212.8) and (1076.5,185.1) .. (1073.7,154.7);
  \draw[line width=5.0pt] (1073.7,154.7) .. controls (1069.1,148.1) and (1061.2,143.0) .. (1053.0,143.0);
  \draw[line width=5.0pt] (1051.0,143.0) .. controls (1046.1,163.9) and (1033.4,182.5) .. (1016.0,196.0);
  \draw[line width=4.0pt] (730.0,236.0) .. controls (734.4,247.4) and (736.3,261.5) .. (746.0,270.0);
  \draw[line width=5.0pt] (836.0,352.0) .. controls (896.6,312.9) and (957.7,268.4) .. (1007.0,213.0);

  %% ════ 圆 / 椭圆（probe 精测）════
  \draw (862.2,272.0) circle (172.3);

  %% ════ 实心点 ════
  \fill (1061.5,112.6) circle (5.6);
  \fill (949.0,176.0) circle (4.0);
  \fill (942.0,399.0) circle (5.1);

  %% ════ 标签（probe 直接给的 \node 行）════
  %% ⚠ 全部补了 fill=white：文字在最上层，否则与曲线交叉干扰
     \node[anchor=center,inner sep=0pt,fill=white,font=\fontsize{31.4}{39.3}\selectfont] at (527.1,134.6) {\textsf{\textbf{\textit{g}}}};   % box(516,121,18,24)
     \node[anchor=center,inner sep=0pt,fill=white,font=\fontsize{33.2}{41.5}\selectfont] at (862.8,273.7) {\textsf{\textbf{\textit{o}}}};   % box(853,263,18,18)
     \node[anchor=center,inner sep=0pt,fill=white,font=\fontsize{30.7}{38.3}\selectfont] at (526.8,376.8) {\textsf{\textbf{\textit{b}}}};   % box(517,364,17,23)
     \node[anchor=center,inner sep=0pt,fill=white,font=\fontsize{21.5}{26.9}\selectfont] at (218.0,477.9) {\textsf{\textbf{1}}};   % box(212,470,7,15)
     \node[anchor=center,inner sep=0pt,fill=white,font=\fontsize{30.9}{38.6}\selectfont] at (202.3,485.5) {\textsf{\textbf{\textit{B}}}};   % box(191,473,19,24)
     \node[anchor=center,inner sep=0pt,fill=white,font=\fontsize{24.9}{31.2}\selectfont] at (900.1,492.4) {\textsf{\textbf{2}}};   % box(893,484,13,16)
     \node[anchor=center,inner sep=0pt,fill=white,font=\fontsize{30.0}{37.5}\selectfont] at (882.0,501.0) {\textsf{\textbf{\textit{R}}}};   % box(872,489,20,22)
     \node[anchor=center,inner sep=0pt,fill=white,font=\fontsize{31.1}{38.8}\selectfont] at (941.3,500.8) {\textsf{\textbf{6}}};   % box(933,489,15,22)
     \node[anchor=center,inner sep=0pt,fill=white,font=\fontsize{30.3}{37.9}\selectfont] at (924.8,506.6) {{\mfallback\textbf{−}}};   % box(909,500,17,4)

  % 钉死画布：否则超出的图元/控制点会把页面撑大，1:1 回比就失准
  \pgfresetboundingbox
  \path[use as bounding box] (0,0) rectangle (1088,522);
\end{tikzpicture}
\end{document}

% ⚠ 待核清单（probe 拿不准，人眼过一遍）：
%   · 可疑字形：\node[anchor=center,inner sep=0pt,font=\fontsize{33.2}{41.5}\selectfont] at (862.8,273.7) {\textsf{\textbf{\te
%   · 未识别/跳过：⑤ 跳过/未识别的字形
